% Focused action-history replication. The broad steering study is a separate paper.
\documentclass{article}
\PassOptionsToPackage{numbers,sort&compress}{natbib}
\usepackage[dblblindworkshop]{neurips_2026}
\workshoptitle{Long-Context Foundation Models}
\usepackage[T1]{fontenc}
\usepackage{microtype}
\usepackage{amsmath,amssymb,booktabs,array,graphicx,float,xcolor}
\usepackage{caption}
\captionsetup{font=small,labelfont=bf}
\usepackage[hidelinks]{hyperref}
\newcommand{\MSE}{\operatorname{MSE}}
\newcommand{\figroot}{../../docs/figures}
% Generated from complete 200-scenario summary SHA256 3769fa71bad0f6e954a6a801a7c23df9ed6abc96ff53958caf5e4a2dc628883e
\newcommand{\LCFMReachChanged}{42}
\newcommand{\LCFMReachCI}{31.6--53.2\%}
\newcommand{\LCFMReachCost}{1.81\%}
\newcommand{\LCFMReachCostCI}{1.24--2.44\%}
\newcommand{\LCFMWallChanged}{59}
\newcommand{\LCFMWallCI}{47.8--69.3\%}
\newcommand{\LCFMWallCost}{1.87\%}
\newcommand{\LCFMWallCostCI}{1.41--2.36\%}
\newcommand{\LCFMCostRange}{1.81--2.56\%}
\newcommand{\LCFMReconRange}{0.4996--0.5117}
\newcommand{\LCFMRankRange}{0.886--0.890}
\newcommand{\LCFMEliteRange}{5.77--5.94}

\title{When Local Interventions Fail as Context Evolves\\in a Predictive World Model}
\author{Anonymous authors}
\begin{document}
\raggedbottom
\maketitle

\begin{abstract}
An internal intervention can reproduce a changed input's immediate prediction and still fail to reproduce its later consequences. We test this in frozen JEPA-WM, which predicts candidate action sequences to support planning. An action appears twice in its rolling two-frame context. Replacing that action's conditioning at its first appearance reproduces the immediate changed-input forecast exactly. When the original action returns as history, the predictions diverge. In a prospective replication on 200 independent starting states, this omission changes the lowest-cost candidate in \LCFMReachChanged/100 Reach states and \LCFMWallChanged/100 Reach-Wall states. The selected candidates incur \LCFMReachCost{} and \LCFMWallCost{} higher mean cost under the changed-input reference. Replacing the conditioning at both appearances restores exact forecast and selection agreement. The experiment establishes a concrete failure of one-step intervention validation: an edit must remain consistent with the information reused during later predictions. It tests a two-frame context and six-step horizon; whether the mismatch grows with longer context remains open.
\end{abstract}

\section{Introduction}
A world model lets a planner compare actions through their predicted consequences. JEPA-WM performs this comparison in learned visual and robot-state features \citep{terver2025jepawm}. An internal intervention intended to represent a changed action should therefore be checked against the predictions and choices produced by changing that action at the input. Agreement at the next prediction is only the first check.

We study a specific reason that agreement can fail later. The predictor reuses each action in two consecutive context windows. A patch at the first appearance leaves the second appearance unchanged. The immediate computation matches the changed-input reference, but the next computation combines the changed predicted state with the original action history. This is a testable mismatch between two computations of the same frozen model.

An exploratory sixteen-state study identified this effect. We fixed the intervention definitions, fresh inputs, and analysis before running a separate 200-state replication. The contribution is a measured connection from local forecast agreement to later candidate selection, together with a control that removes the mismatch. We use ``counterfactual'' to mean the model's response to a changed input action; physical counterfactual accuracy is not measured.

\begin{figure}[t]
\centering
\includegraphics[width=\textwidth]{\figroot/lcfm_context_lifetime.pdf}
\caption{\textbf{A locally exact action patch can change the later choice.} (A) Action-conditioning occurrences in successive windows. The changed action $a_3^*$ occupies the newest position at H3 and the older position at H4. (B) Mean distance to the changed-input forecast, divided by the unmodified forecast's distance (Equation~\ref{eq:distance}); zero is exact agreement. Bands are marginal 95\% state-bootstrap intervals. (C) H6 candidate disagreement in the original bank, with 97.5\% Wilson intervals per task. Mean extra reference cost includes states with unchanged choices. Both tasks have 100 independent states. Patching both appearances gives exact agreement at all measured horizons.}
\label{fig:lifetime}
\end{figure}

\newpage
\section{Experiment}
\label{sec:experiment}
\paragraph{Model and reference.}
We use the released MetaWorld JEPA-WM checkpoint, frozen in strict FP32 with TF32 disabled. Its 39.69M parameters include a 22.06M visual encoder, a 17.63M predictor, and a small proprioceptive embedding. The six-block predictor receives two frames of context and forecasts six imagined steps. For each starting state, we retain two banks of 300 candidate action sequences. The first bank is the primary evaluation; the second checks sensitivity to candidate sampling on the same states.

For candidate $i$, the reference replaces its H3 action with candidate $(i+1)\bmod 300$'s H3 action and rolls out the model normally. All other input actions stay fixed. Each model action groups five elementary controls. This cyclic substitution supplies a concrete input intervention without optimizing a donor for a favorable outcome.

\paragraph{Internal replacements.}
We record the action-conditioning tensors from the changed-input computation. The \emph{one-time patch} replaces the H3 action's conditioning at all six blocks only when that action first appears. The \emph{persistent patch} also replaces its older-position occurrence at H4. Predicted states then evolve within each patched rollout. The raw action array remains unmodified in the patched runs. A persistent all-block patch reproduces every conditioning input changed by the reference, so exact agreement is the expected consistency control. It is not a learned correction.

The fixed registry contains 35 conditions: unmodified, capture-only, changed-input reference, the two all-block patches, and five conditions at each of six individual blocks. These single-block conditions are one-time replacement, persistent replacement, and three norm-matched random controls in the input-induced range, its orthogonal complement, and the full conditioning space. Appendix~\ref{app:controls} retains every block and control.

\paragraph{Measurements.}
Let $\hat z^{P}_{i,h}$, $\hat z^{CF}_{i,h}$, and $\hat z^{N}_{i,h}$ be patch, changed-input, and unmodified predictions. For a state and candidate bank, define
\begin{equation}
 \begin{aligned}
 D_h(P,CF)&=\frac{1}{300}\sum_i\left[\MSE(\hat z^{P,v}_{i,h},\hat z^{CF,v}_{i,h})+0.1\MSE(\hat z^{P,p}_{i,h},\hat z^{CF,p}_{i,h})\right],\\
 E_h(P)&=\frac{D_h(P,CF)}{D_h(N,CF)}.
 \end{aligned}
 \label{eq:distance}
\end{equation}
The ratio is calculated within each state before averaging. $E=0$ means exact agreement; $E=1$ is the unmodified forecast's discrepancy. We report H3, H4, and H6. The accompanying reconstruction score is $R=1-E$.

The planner's terminal cost is $C_i=\MSE(\hat z^v_{i,6},z_g^v)+0.1\MSE(\hat z^p_{i,6},z_g^p)$. The primary endpoint is disagreement between the patch's and reference's selected minimum-cost candidate at H6. To measure its size, we evaluate the patch-selected index $j_P$ using reference costs:
\begin{equation}
 G(P)=100\,\frac{C^{CF}_{j_P}-\min_i C^{CF}_i}{\min_i C^{CF}_i}.
 \label{eq:cost}
\end{equation}
This is excess cost according to the reference model. We also retain Spearman rank agreement and overlap among the ten lowest-cost candidates. These fixed-bank comparisons do not execute the selected actions or continue adaptive search.

\paragraph{Replication and inference.}
We sample 100 independent starting states each for Reach and Reach-Wall, excluding the earlier sixteen states and engineering inputs. All 35 conditions and both banks for a state run on the same GPU. The workload is 14,000 batched six-step forecasts, or 4.2 million candidate forecasts. Candidates, banks, and intervention arms are repeated measurements within a state.

Each primary task rate has a 97.5\% Wilson interval, using a Bonferroni correction across the two rates for nominal 95\% family coverage. Secondary means use 20,000 paired state-bootstrap draws with shared resampling weights across arms and banks. These intervals are marginal and descriptive. All registered cells are retained, with no layer selection or pooling of the exploratory cohort. Appendix~\ref{app:repro} gives the evidence identities and execution details.

\section{Results}
\paragraph{The mismatch begins when the action returns as history.}
Both all-block patches match the changed-input forecast exactly at H3. At H4, the one-time patch's mean normalized discrepancy rises to approximately 0.53 on both tasks; it remains about 0.49--0.50 at H6. The persistent patch has zero discrepancy at every measured horizon in both banks (Figure~\ref{fig:lifetime}). Thus the mismatch appears at the omitted contextual occurrence. It does not grow monotonically from H4 to H6.

\paragraph{The mismatch reaches candidate selection.}
The one-time patch changes the selected candidate in \LCFMReachChanged/100 Reach states (interval \LCFMReachCI) and \LCFMWallChanged/100 Reach-Wall states (\LCFMWallCI). Mean excess reference cost is \LCFMReachCost{} on Reach (95\% interval \LCFMReachCostCI) and \LCFMWallCost{} on Reach-Wall (\LCFMWallCostCI). The cost differences are modest even though a substantial fraction of choices change. Across both tasks and banks, mean Spearman agreement is \LCFMRankRange{} and mean top-ten overlap is \LCFMEliteRange{} candidates. Replacing both appearances restores exact ranking and winner agreement, with zero excess reference cost.

\paragraph{A second bank gives a similar forecast discrepancy.}
Across the four task/bank cells, the one-time patch's H6 reconstruction score is \LCFMReconRange{} and mean excess cost is \LCFMCostRange{}. The second bank tests sampling sensitivity within the same 200 states. The complete single-block and random-control grid is included in Appendix~\ref{app:controls}; none of those comparisons changes the primary all-block estimate.

\newpage
\section{Interpretation and scope}
\label{sec:scope}
A one-step patch check can accept an intervention whose later rollout no longer represents its intended input change. Here the cause is directly inspectable: the unchanged action is reused in the next window. Replacing that occurrence removes the mismatch. Exact persistent agreement follows from covering the changed conditioning inputs, while the experiment measures how much omitting one occurrence matters for forecast and selection.

This differs from ordinary forecast error against the environment. Both sides use the same checkpoint, starting state, and candidate bank, so the comparison isolates intervention consistency within the model. It establishes neither physical accuracy of the changed-input forecast nor improved robot success. Selection frequency also depends on the candidate cost gaps; many changed winners can have similar costs, as the small mean cost differences demonstrate.

Activation-patching studies show that the metric and intervention definition affect conclusions \citep{zhang2023patching}. Physical representation studies also use interventions to test what a video model encodes \citep{joseph2026interpretingphysics}. Our experiment adds a temporal check for an action-conditioned world model: validate the intervention after the relevant input is reused, through the forecast horizon that supplies the planning decision.

The present evidence covers one checkpoint, two tasks, and a two-frame context. LeWorldModel provides a possible second-model test with three-frame history on PushT \citep{maes2026lewm}; it would extend the number of explicit action occurrences by one. Establishing context-length scaling would require matched models trained for different history lengths and separate control of rollout horizon. More observations from the present checkpoint would narrow its uncertainty without answering that question.

\clearpage
\bibliographystyle{plainnat}
\bibliography{references}
\clearpage
\appendix
\section{Complete forecast and selection measurements}
\begin{figure}[H]
\centering\includegraphics[width=\textwidth]{\figroot/lcfm_replication_history.pdf}
\caption{Both candidate banks on the same 100 states per task. Reconstruction is $R=1-E$ from Equation~\ref{eq:distance}. Means and marginal 95\% paired state-bootstrap intervals retain every state. Coincident persistent-patch curves are exact reference agreement.}
\end{figure}
\begin{figure}[H]
\centering\includegraphics[width=\textwidth]{\figroot/lcfm_replication_choices.pdf}
\caption{Primary selection disagreement and secondary reference cost. The second bank is a repeated measurement of the same states. All cost means include unchanged selections, whose excess cost is zero.}
\end{figure}
\clearpage
\section{All blocks and random controls}
\label{app:controls}
\begin{figure}[H]
\centering\includegraphics[width=\textwidth]{\figroot/lcfm_replication_layers.pdf}
\caption{Complete six-block by five-condition grid, separately for both tasks and banks. Values are mean H6 reconstruction $R=1-E$. Negative values mean greater forecast discrepancy than the unmodified model. Every cell contains 100 states; all 5,460 registered secondary summary cells across metrics and horizons are defined. The color scale is shared and includes zero.}
\end{figure}
Single-block replacement leaves the other blocks' conditioning unchanged and need not reproduce the reference. Random controls match the prescribed per-row norm and retain their fixed seeds. These comparisons are descriptive checks of intervention scope and direction. The all-block result uses no layer chosen from this grid.

\section{Reproduction and evidence}
\label{app:repro}
The execution protocol was frozen before the fresh replication and records all 200 state seeds, two candidate banks, 35 conditions, the checkpoint, and precision. SHA256 identities (line-wrapped for readability) are:
\begin{center}\small
\begin{tabular}{ll}
\toprule
Protocol & \texttt{967aebfd743fe2bfb226bef2915dbebbac3}\\
 & \texttt{1f4c816f84f987ed514a3356c9204}\\
Checkpoint & \texttt{c3297772c7af4e84c28bd0c0937e22398a8}\\
 & \texttt{888485305014bb3ab3f1508a310bd8}\\
Summary & \texttt{3769fa71bad0f6e954a6a801a7c23df9ed}\\
 & \texttt{6abc96ff53958caf5e4a2dc628883e}\\
\bottomrule
\end{tabular}
\end{center}
The analysis records 42,000 ranking rows and 126,000 forecast-reconstruction rows. The primary count compares the deterministic selected index. Tied-minimum membership is retained separately; Spearman uses average ranks for ties. Undefined ratios or constant rank vectors remain explicit rather than being dropped; all registered secondary cells in this replication are defined. The bootstrap seed is 2026091421. State, goal, and environment-seed exposure checks establish separation from audited project records, without claiming novelty relative to base-model pretraining.

The public analysis table, row-level compressed exports, number generator, and vector-figure generator share a hash-checked summary. A separate checker recounts the primary events, recomputes Wilson intervals with SciPy, and checks every secondary mean. The 200 completed raw records and derived outputs were preserved with verified cloud readback before releasing the workers. Identifying repository and storage links are omitted here; this draft does not yet include an anonymous raw-data package.

\paragraph{Compute.}
Whole-state replicas keep all paired conditions together, batching 300 candidates per forward call. Productive workers used RTX 4090 and RTX 5090 GPUs, with a peak fleet of two 4090s and eight 5090s. Representative complete states took 223--230 seconds on a 4090 and 185.5--189.8 seconds on a 5090. The measured peak allocated memory was 18.38\,GB. Raw compressed scientific archives total 1.576\,GB; execution retained strict FP32, TF32 disabled, and pinned inputs. Twelve engineering case executions and three interrupted technical partials were excluded from the scientific cohort and preserved separately. The earlier exploratory study and broader steering project consumed additional compute; the figures here use only the complete replication.

\paragraph{Broader impact.}
Checking whether an internal intervention reproduces its intended input change can prevent misleading interpretations of world-model behavior. Agreement with a model-generated reference does not validate physical safety. This experiment involves simulator starting states and model forecasts; it introduces no real-robot deployment or human-participant study.
\clearpage
\input{checklist.tex}
\end{document}
